Nach dem Stefan-Boltzmann-Gesetz ist die Strahlungsintensität $J = \varepsilon\,\sigma\,T^4$ proportional zur vierten Potenz der \emph{absoluten} Temperatur. Mit $T_1 = \Tv$ und $T_2 = \Tn$ ist
\begin{align}
 \frac{J_2}{J_1} &= \left(\frac{T_2}{T_1}\right)^4
\end{align}
und die relative Zunahme
\begin{align}
 \frac{J_2-J_1}{J_1} &= \incF \\
   &= \left(\frac{\Tn}{\Tv}\right)^4-1 \\
   &= \inc \approx \result{\incP{0}{2}}
\end{align}
